\note{sa page 11}

\[ (6.15) \qquad \begin{cases} L \Lambda = 1 - \sum_{0 \leqslant i \leqslant n} \pi_{i,0} \\[2pt] \Lambda L = 1 - \sum_{n \leqslant j \leqslant 2n} \pi_{j,0} \end{cases} \]

\struck{Quant aux $\pi_{j,0}$, ils sont déterminés par les \ill{}}

\marginal{\uncertain{cf n° 7} pour l'anneau $\Phi$ : \uncertain{syst.\ de} relations \ldots}

\emph{Question.} Donner une \ill{} synthétique des relations entre les
$\pi_{i}$ comme \struck{\ill{}} $\Lambda$, $L$ exclusivement, et \ill{} des
relations analogues \ill{} $L$, $\Lambda$ \ill{} qui définissent $\Lambda$,
\uncertain{polynômes non commutatifs} en $L$ et des $\pi_{i}$ \ldots

7. \uncertain{Revenant} à un anneau gradué abstrait $\mathcal{E}$, muni d'un
$L \in \mathcal{E}^{(2)}$, on veut \uncertain{obtenir} la situation
\uncertain{du n° précédent}, sans utiliser $M^{*}$, en termes de propriétés du
seul anneau gradué $\mathcal{E}^{*}$ muni de $L \in \mathcal{E}^{(2)}$.
Partons pour ceci \struck{\ill{}} d'un module libre $M_{0}$ de rang $2n$ sur
$\mathbb{Z}$, muni d'une $2$-forme alternée \uncertain{non dégénérée}
\add{(sur $\mathbb{Z}$ !)} $\xi_{0} \in \Lambda^{2} M_{0}$, \add{$\xi^{*}_{0} \in
\Lambda^{2} M^{\vee}_{0}$ la forme duale,} soit $M^{*}_{0} = \Lambda M_{0}$,
% the index of M_0 and xi_0 is written like an « a » on this page; End(M^*_0) and Phi_0 below fix it as 0.
et considérons sur $M^{*}_{0}$ les opérateurs
\[ L_{0}\, x = \xi_{0}\, x , \qquad \Lambda_{0}\, x = \xi^{*}_{0} \lrcorner\, x \]
et considérons l'anneau \add{gradué}
\[ \Phi_{0} = \mathbb{Z}[L_{0}, \Lambda_{0}] \subset \operatorname{End}(M^{*}_{0}) . \]
Notons que $\operatorname{End}_{\mathbb{Z}}(M^{*}_{0})$ \add{opérant sur
$M^{*}_{\mathbb{Z}}$}, muni de $L_{0}$, satisfait à la condition
de $(6.1)$\uncertain{$(\omega)$}, \struck{\ill{}} et que $\Lambda_{0}$ est
l'opérateur associé défini par (6.15).

\note{Le $\mathbb{Z}$ est porté sous $\operatorname{End}$, et « opérant sur
$M^{*}_{\mathbb{Z}}$ » est une insertion au-dessus de la ligne, rattachée par un
long trait. Ce que la parenthèse après $(6.1)$ contient est un seul glyphe
arrondi, lu $\omega$ et non identifié autrement. La phrase identifie les deux
$\Lambda_{0}$ du feuillet : celui que la formule affichée définit par la
contraction $\xi^{*}_{0} \lrcorner\,-$, et celui que caractérisent les relations
(6.15).}

